% Section 6.2, from lectures/L06/appendix/b_mmio_and_resources.md.

\section{MMIO and Resource Management}
\label{c6:app:b}

How a driver reaches its hardware: finding the address, claiming it, mapping it, and reading it
without the compiler or the processor helpfully rearranging what you asked for.

Every listing here was produced on this course's target, against the \code{qa-dev} device. Its
register map is \file{tools/qa-dev/qa-dev.h}, and that header is the only place the map exists: the
QEMU device model includes it, and so does every driver from here on, through a copy in its lab
directory. Read it as you would a datasheet. Here it is, typeset straight from the repository:

\cfile{tools/qa-dev/qa-dev.h}

\subsection{Where the address comes from}
\label{c6:app:b1}

On an embedded system, nothing discovers it. Somebody wrote it down, and in a device tree that
somebody is the board's author. \code{qa-dev}'s node, as QEMU generated it:

\begin{console}
platform-bus@c000000 {
        ranges = <0x00 0x00 0xc000000 0x2000000>;
        #address-cells = <0x01>;
        #size-cells = <0x01>;

        qa-dev@0 {
                compatible = "qacademy,qa-dev-1.0";
                reg = <0x00 0x1000>;
                interrupts = <0x00 0x70 0x04>;
        };
};
\end{console}

\noindent \textbf{\code{reg = <0x00 0x1000>} does not say the device is at address zero.} It says
the device begins at offset 0 \emph{on its parent bus}, and is 0x1000 bytes long. To get a physical
address you translate through the parent's \code{ranges}, which reads: child 0x0 maps to parent
0xC000000, for 0x2000000 bytes. So:

\[
  \texttt{0x0} + \texttt{0xC000000} = \textbf{\texttt{0xC000000}}
\]

\noindent A reader who takes \code{reg} at face value maps address 0 and reads nothing useful. That
translation is the Cross-check, and Chapter~\ref{c10:ch} is where \code{platform_get_resource} does
it for you.

\textbf{The cells are 32-bit big-endian}, whatever the machine's own endianness. On the target the
properties are raw bytes:

\begin{shell}
od -An -tx1 -v /sys/firmware/devicetree/base/platform-bus@c000000/qa-dev@0/reg
\end{shell}

\noindent Reading them with \code{od -tx4} is the obvious shortcut and is wrong: it decodes four
bytes as a \emph{native-endian} integer, so on this little-endian target every cell comes back
byte-swapped and \code{0x0C000000} reads as \code{0x0000000C}. That number is plausible enough to
act on, which is what makes it worth knowing. The lab's \file{run.sh} reassembles bytes by hand for
exactly this reason.

\lecturefigure{lectures/L06/appendix/images/address_translation.png}
  {Four boxes in a row showing one address being translated: reg holding 0x00 and 0x1000, ranges
   mapping child 0x0 to 0x0C000000, ioremap taking that physical address, and readl returning
   0x51414456. Below, a note that a wrong mapping fails two ways and only one of them is loud.}
  {fig:address_translation}

\subsection{Claiming the region}
\label{c6:app:b2}

\begin{ccode}
region = request_mem_region(base, QA_DEV_MMIO_SIZE, "qa_mmio");
if (!region)
        return -EBUSY;
\end{ccode}

\noindent \textbf{This maps nothing and prevents nothing.} It is an entry in a registry. What it
buys you is that a \emph{second} driver calling it for an overlapping range gets NULL, so two
drivers cannot both believe they own the same device and work intermittently.

It is bookkeeping that the hardware knows nothing about. A driver that skips it and calls
\code{ioremap} directly works fine, right up until the day a second driver does the same.

The registry is visible, and this is the payoff for Exercise~\ref{c1:ex:7}, which looked for
\code{qa-dev} in \file{/proc/iomem} and found nothing:

\begin{console}
# grep qa_mmio /proc/iomem
0c000000-0c000fff : qa_mmio
\end{console}

\noindent The name in that listing is the string you passed, so make it the driver's name and not
something generic. Release with \code{release_mem_region(base, size)}; failing to means the next
\code{insmod} gets \code{-EBUSY} from your own previous load, which is a confusing way to discover a
leak.

\subsection{Mapping it}
\label{c6:app:b3}

\begin{ccode}
regs = ioremap(base, QA_DEV_MMIO_SIZE);
if (!regs)
        return -ENOMEM;
\end{ccode}

\noindent \code{ioremap} creates a kernel virtual mapping for a physical range and returns
\code{void __iomem *}.

\textbf{It is not memory, and the distinction is not pedantry.} Four things are true of device
registers and false of memory:
\begin{itemize}
  \item \textbf{Reads have side effects.} Reading \code{QA_DEV_SAMPLE} removes a sample from the
    FIFO. Reading it twice is not the same as reading it once.
  \item \textbf{Writes may not be merged or reordered.} Two writes to the same register are two
    operations.
  \item \textbf{It must not be cached.} \code{ioremap} maps it uncached; a cached mapping would let
    the processor satisfy a read from a stale line and postpone a write indefinitely.
  \item \textbf{There is no storage.} Nothing was allocated. \code{iounmap} returns the mapping, not
    any memory.
\end{itemize}
Which is why the compiler must not be trusted with it. Given an ordinary pointer, a compiler may
eliminate a read whose value you ignore, merge two writes, or hoist a read out of a loop. All three
are correct for memory and all three break a driver.

\subsection{\code{readl} and \code{writel}}
\label{c6:app:b4}

\begin{ccode}
u32 id = readl(regs + QA_DEV_ID);
writel(0xdeadbeef, regs + QA_DEV_SCRATCH);
\end{ccode}

\noindent These are the only legal way to touch an \code{__iomem} pointer. They compile to a single
load or store plus the barriers the architecture needs, and they are not optimised away.

\begin{narrowtable}{@{}ll@{}}
  \thead{Accessor} & \thead{Width} \\
  \midrule
  \code{readb}/\code{writeb} & 8 bits \\
  \code{readw}/\code{writew} & 16 bits \\
  \code{readl}/\code{writel} & 32 bits \\
  \code{readq}/\code{writeq} & 64 bits \\
\end{narrowtable}

\noindent \textbf{Use the width the device documents.} \code{qa-dev} accepts 32-bit accesses only
and rejects anything else rather than quietly widening it, which is deliberate: a driver that gets
away with a byte write to a word register here would not get away with it on real hardware.

The address arithmetic is plain byte arithmetic on the returned pointer, so \code{regs + QA_DEV_ID}
with the offsets from the register map is exactly right, and casting to a struct of \code{u32}
fields is a common alternative that works until a register map has a gap in it.

\subsection{Ordering, and the \code{_relaxed} variants}
\label{c6:app:b5}

\code{writel} carries a barrier; \code{writel_relaxed} does not. What the barrier orders is the part
worth getting right.

Accesses from one CPU to one device's registers, through a mapping \code{ioremap} made, arrive at
the device in program order either way. Write a control register and then read the status
register, relaxed or not, and the device sees the write first. What the barrier orders is a register
access against \emph{ordinary memory}: \code{writel} waits for every earlier store to RAM before its
own store, and \code{readl} completes before any later load from RAM begins. That matters the moment
the device reads memory itself. Fill a descriptor in RAM, then write the register that tells the
device to fetch it:

\begin{ccode}
desc->len = len;                          /* ordinary memory */
writel(DOORBELL_GO, regs + DOORBELL);     /* ordered after it */
\end{ccode}

\noindent Without the barrier the processor may let the register write overtake the store to the
descriptor, and the device fetches a descriptor from before your own update. \code{qa-dev} reads no
memory, so every access this course's labs make is to one device, and the difference never shows
on this target.

The relaxed forms exist because barriers cost time, and a driver writing a burst of
registers where only the last one matters can use \code{_relaxed} for the burst and an ordered
access at the end. \textbf{Use the ordered forms until you have measured a reason not to.} A missing
barrier produces a driver that works on one machine and fails on another, which is the single worst
class of bug to debug.

This is the same memory ordering that is a property of the processor on its own; here there is
hardware on the other end that also has an opinion.

\subsection{What a wrong address does}
\label{c6:app:b6}

Worth knowing before you guess one, and the answer is not ``you read a wrong number''.

\textbf{A wrong-but-mapped address reads plausibly.} Pointing the lab's module at \code{0xC} rather
than \code{0xC000000} produces:

\begin{console}
qa_mmio: id=0x00000000 version=0x0000
\end{console}

\noindent No error, no fault, just zeros. A driver that did not check its identity register would
carry on and misbehave later.

\textbf{A wrong-and-unmapped address aborts the machine.} Pointing it one page past the device, at
\code{0xC001000}, produces this:

\begin{console}
Internal error: synchronous external abort: 0000000096000010 [#1] PREEMPT SMP
Modules linked in: qa_mmio(O+)
CPU: 0 UID: 0 PID: 67 Comm: insmod Tainted: G           O       6.12.30
pc : qa_mmio_init+0x98/0xff8 [qa_mmio]
Call trace:
 qa_mmio_init+0x98/0xff8 [qa_mmio]
 do_one_initcall+0x84/0x360
\end{console}

\noindent \code{insmod} is killed, the module is left half-loaded so it can be neither used nor
removed, and the sensible next step is a reboot. On this target that costs two seconds; on a board
on a bench it costs a walk across the room, and on a product in the field it costs a truck.

Three lessons, and the third is the one worth carrying:
\begin{itemize}
  \item \textbf{Always check an identity register} if the device has one. \code{qa-dev} has
    \code{QA_DEV_ID} for this, and reading \code{0x51414456} is what tells you the mapping landed
    where you meant.
  \item \textbf{The failure mode depends on the address}, so ``it did not crash'' is not evidence
    that a mapping is right.
  \item \textbf{This is what the emulator is for.} \Secref{c1:app:a4} said the course cannot teach
    you what a bricked board feels like; the other side of that trade is that you can point a driver
    at a wrong address here and find out what happens, which is not an experiment anyone runs twice
    on real hardware.
\end{itemize}

\subsection{The whole sequence, and its unwind}
\label{c6:app:b7}

Acquire in order, release in reverse, and unwind only what you actually got:

\begin{ccode}
region = request_mem_region(base, QA_DEV_MMIO_SIZE, "qa_mmio");
if (!region)
        return -EBUSY;

regs = ioremap(base, QA_DEV_MMIO_SIZE);
if (!regs) {
        ret = -ENOMEM;
        goto err_region;
}

if (readl(regs + QA_DEV_ID) != QA_DEV_ID_MAGIC) {
        ret = -ENODEV;
        goto err_map;
}
return 0;

err_map:
        iounmap(regs);
err_region:
        release_mem_region(base, QA_DEV_MMIO_SIZE);
        return ret;
\end{ccode}

\noindent Two acquisitions, two labels, in reverse. Add a third and you add a third label; the
pattern does not get cleverer, it gets longer, and by the time a real driver has a clock, a
regulator, a reset line and three mappings it is twenty lines of unwinding that must stay in step
with the acquire path.

That is the argument for \code{devm_}, which \S\ref{c6:app:a7} explains and which this chapter
cannot use, because every \code{devm_} function needs a \code{struct device *} and a bare module
does not have one. Chapter~\ref{c10:ch} is where one appears and this ladder is deleted.
